% !TeX program = lualatex
\documentclass[aspectratio=169,t]{beamer}
\usepackage{../../styles/default}
\usepackage{booktabs}
\usepackage{tabularx}
\usepackage{fontawesome5}
\usetikzlibrary{arrows.meta,positioning,shadows}

% Figure palette shared with /home/beta/paper-code-retrieval/deck.tex.
\definecolor{agentmain}{HTML}{1E40AF}
\definecolor{agentlight}{HTML}{F0F6FF}
\definecolor{agentborder}{HTML}{93C5FD}
\definecolor{successmain}{HTML}{16A34A}
\definecolor{successlight}{HTML}{F0FDF4}
\definecolor{successborder}{HTML}{86EFAC}
\definecolor{failmain}{HTML}{DC2626}
\definecolor{faillight}{HTML}{FEF2F2}
\definecolor{failborder}{HTML}{FCA5A5}
\definecolor{couragemain}{HTML}{C2410C}
\definecolor{couragelight}{HTML}{FFF7ED}
\definecolor{courageborder}{HTML}{FDBA74}
\definecolor{textmid}{HTML}{4B5563}
\definecolor{textlight}{HTML}{94A3B8}

% The five benchmark tiers reuse the semantic palette.
\colorlet{tierone}{successmain}
\colorlet{tiertwo}{couragemain}
\colorlet{tierthree}{agentmain}
\colorlet{tierfour}{failmain}
\colorlet{tierfive}{textlight}
\colorlet{accentblue}{agentmain}
\colorlet{accentred}{failmain}
\colorlet{accentgreen}{successmain}

% Every palette exposes <base>main / <base>light / <base>border so one figure
% macro can take the palette name as an argument.
\colorlet{neutralmain}{textmid}
\colorlet{neutrallight}{blackPrimary!2}
\colorlet{neutralborder}{blackPrimary!18}

\tikzset{
  figure card/.style={rounded corners=3pt,line width=0.6pt,
  drop shadow={opacity=0.08,shadow xshift=0.4mm,shadow yshift=-0.4mm}},
  % Card text stays dark for legibility; the palette colors the icon and label.
  agent card/.style={figure card,draw=agentborder,fill=agentlight,text=blackPrimary},
  success card/.style={figure card,draw=successborder,fill=successlight,text=blackPrimary},
  fail card/.style={figure card,draw=failborder,fill=faillight,text=blackPrimary},
  courage card/.style={figure card,draw=courageborder,fill=couragelight,text=blackPrimary},
  neutral card/.style={figure card,draw=neutralborder,fill=neutrallight,text=blackPrimary},
  figure arrow/.style={-{Stealth[round,length=2.2mm]},line width=0.9pt,draw=blackPrimary!35},
}

\setbeamercolor{callout}{fg=blackPrimary,bg=blackPrimary!5}
\setbeamercolor{softblock}{fg=blackPrimary,bg=blackPrimary!3}

\newcommand{\callout}[1]{%
  \vspace{0.45em}%
  \begin{beamercolorbox}[wd=\linewidth,sep=0.55em,rounded=false]{callout}
    #1
  \end{beamercolorbox}%
  \vspace{0.35em}%
}

\newcommand{\metric}[2]{%
  \begin{minipage}[t]{\linewidth}\raggedright
    {\fontsize{18}{20}\selectfont\bfseries #1\par}
    \vspace{0.2em}
    {\scriptsize\color{blackSecondary}#2\par}
  \end{minipage}%
}

% -----------------------------------------------------------------------------
% Stage row: n equal-width, equal-height cards spanning the whole text width.
%
%   \begin{stagerow}{<n>}{<card height>}
%     \stage{<palette>}{<icon>}{<label>}{<detail>}   % repeated n times
%   \end{stagerow}
%
% Each card holds a fixed-size minipage, so the row can never go ragged no
% matter how many lines a card's text takes. Cards are named stage1..stagen,
% so extra tikz code (e.g. a feedback arrow) can be added inside the row.
% Pass [plain] for a row of peers rather than a left-to-right pipeline.
% -----------------------------------------------------------------------------
\newif\ifstagearrows
\newcounter{stagenum}
\newcounter{stageprev}
\newlength{\stagegap}\setlength{\stagegap}{0.6cm}
\newlength{\stagew}\newlength{\stageh}\newlength{\stagex}
\newlength{\stageinnerw}\newlength{\stageinnerh}
\newlength{\stagepad}\setlength{\stagepad}{0.5em}

\tikzset{
  stage box/.style={inner xsep=\stagepad,inner ysep=\stagepad,font=\small},
}

\newenvironment{stagerow}[3][arrows]{%
  \stagearrowstrue
  \ifstrequal{#1}{plain}{\stagearrowsfalse}{}%
  \setcounter{stagenum}{0}%
  \setlength{\stagew}{\dimexpr(\linewidth-\numexpr#2-1\relax\stagegap)/#2\relax}%
  \setlength{\stageh}{#3}%
  \setlength{\stageinnerw}{\dimexpr\stagew-2\stagepad-1.2pt\relax}%
  \setlength{\stageinnerh}{\dimexpr\stageh-2\stagepad-1.2pt\relax}%
  \noindent
  \begin{tikzpicture}
  }{%
  \end{tikzpicture}%
}

% Cards sit at computed offsets rather than relative to their predecessor, so
% the row is a fixed grid: colored icon, dark label, gray detail.
\newcommand{\stage}[4]{%
  \setcounter{stageprev}{\value{stagenum}}%
  \stepcounter{stagenum}%
  \setlength{\stagex}{\dimexpr\numexpr\value{stagenum}-1\relax\dimexpr\stagew+\stagegap\relax\relax}%
  \node[stage box,#1 card,anchor=north west] (stage\thestagenum) at (\the\stagex,0)
  {
    \begin{minipage}[t][\stageinnerh][t]{\stageinnerw}\centering%
      \ifblank{#2}{}{{\color{#1main}#2}\par\vspace{0.3em}}%
      {\bfseries #3\par}%
      \ifblank{#4}{}{\vspace{0.2em}{\scriptsize\color{blackSecondary}#4\par}}%
  \end{minipage}};
  \ifstagearrows
  \ifnum\value{stagenum}>1
  \draw[figure arrow] (stage\thestageprev) -- (stage\thestagenum);
  \fi
  \fi
}

% -----------------------------------------------------------------------------
% Side-by-side panels of equal height (used for the two-column contrast slides).
%   \panel{<palette>}{<height>}{<heading>}{<body>}
% -----------------------------------------------------------------------------
\newcommand{\panel}[4]{%
  \begin{tikzpicture}
    \node[figure card,draw=#1border,fill=#1light,text=blackPrimary,%
    inner xsep=0.7em,inner ysep=0.65em]%
    {
      \begin{minipage}[t][#2][t]{\dimexpr\linewidth-1.4em-1.2pt\relax}%
        {\bfseries\color{#1main}#3\par}%
        \vspace{0.45em}%
        #4%
    \end{minipage}};
  \end{tikzpicture}%
}

% Tight itemize for use inside \panel and table cells.
\newenvironment{panelitems}{%
\begin{itemize}\small\setlength{\itemsep}{0.3em}%
    }{%
  \end{itemize}%
}

\title{ExploitBench: A Capability Ladder Benchmark for LLM Cybersecurity Agents}
\subtitle{Paper presentation --- Lee \& Brumley (2026), arXiv:2605.14153}
\author{Minjae Gwon}
\institute{POSTECH}
\date{\today}
\brandlogo{figure-postech-wordmark.pdf}
\brandlogoheight{0.0098\paperwidth}
\brandlogovadjust{0pt}

\begin{document}

\begin{frame}
  \titlepage
\end{frame}

\AtBeginSection[]{%
  \sectionframe{currentsection,sectionstyle=show/shaded,subsectionstyle=show/show/hide}%
}

\tocframe{sections={1-3},sectionstyle=show,subsectionstyle=show/show/show}
\tocframe{sections={4-6},sectionstyle=show,subsectionstyle=show/show/show}

% =============================================================================
\section{Preliminaries}
\subsection{Web Pages, Browsers, and V8}

\begin{frame}{What Is a Web Page?}
  \blackfiguretop
  \noindent\makebox[\linewidth][c]{%
    \begin{tikzpicture}
      \path[use as bounding box] (0,0.15) rectangle (13.9,-2.0);

      \node[stage box,neutral card,text width=2.25cm,minimum height=1.45cm,
      align=center] (server) at (1.35,-1.05)
      {{\color{neutralmain}\faServer}\\[0.3em]
      \textbf{Web server}};

      \draw[rounded corners=4pt,draw=blackPrimary!20,line width=0.6pt]
      (3.65,0.02) rectangle (13.45,-1.98);
      \node[fill=white,inner xsep=0.45em,text=blackSecondary,font=\scriptsize\bfseries]
      at (8.55,0.02) {PAGE RESOURCES};

      \node[stage box,agent card,text width=2.55cm,minimum height=1.35cm,
      align=center] (html) at (5.35,-1.05)
      {{\color{agentmain}\faFileCode}\\[0.3em]
        \textbf{HTML}\\[0.2em]
      {\scriptsize\color{blackSecondary}document structure}};
      \node[stage box,courage card,text width=2.55cm,minimum height=1.35cm,
      align=center] (css) at (8.55,-1.05)
      {{\color{couragemain}\faPaintBrush}\\[0.3em]
        \textbf{CSS}\\[0.2em]
      {\scriptsize\color{blackSecondary}visual appearance}};
      \node[stage box,success card,text width=2.55cm,minimum height=1.35cm,
      align=center] (js) at (11.75,-1.05)
      {{\color{successmain}\faCode}\\[0.3em]
        \textbf{JavaScript}\\[0.2em]
      {\scriptsize\color{blackSecondary}interactive behavior}};

      \draw[figure arrow] (server.east) -- node[above,font=\scriptsize,
      text=blackSecondary] {delivers} (3.65,-1.05);
    \end{tikzpicture}%
  }

  \vspace{0.8em}
  \begin{itemize}
    \item A \textbf{web page} is a document and a collection of resources loaded from a site.
      \begin{itemize}
        \item HTML defines headings, buttons, forms, and other structure.
        \item CSS controls layout, color, and typography.
        \item JavaScript reacts to clicks, changes content, and communicates with servers.
      \end{itemize}
    \item Images, fonts, videos, and WebAssembly modules may be loaded alongside those files.
    \item The page's code comes from the website, so a browser must treat it as \textbf{untrusted input}.
  \end{itemize}
\end{frame}

\begin{frame}{What Is a Web Browser?}
  \blackfiguretop
  \noindent\makebox[\linewidth][c]{%
    \begin{tikzpicture}
      \path[use as bounding box] (0,0.30) rectangle (13.9,-2.08);

      \draw[rounded corners=4pt,draw=successborder,line width=0.8pt]
      (0.15,0.18) rectangle (13.75,-2.00);
      \node[fill=white,inner xsep=0.5em,text=successmain,font=\scriptsize\bfseries]
      at (6.95,0.18) {\faLock\quad ISOLATED / SANDBOXED};

      \node[stage box,neutral card,text width=3.15cm,minimum height=1.25cm,
      align=center] (fetch) at (2.2,-0.88)
      {{\color{neutralmain}\faGlobe}\\[0.3em]
        \textbf{Fetch}\\[0.2em]
      {\scriptsize\color{blackSecondary}page files from the network}};
      \node[stage box,agent card,text width=3.15cm,minimum height=1.25cm,
      align=center] (render) at (6.95,-0.88)
      {{\color{agentmain}\faDesktop}\\[0.3em]
        \textbf{Render}\\[0.2em]
      {\scriptsize\color{blackSecondary}HTML and CSS into pixels}};
      \node[stage box,courage card,text width=3.15cm,minimum height=1.25cm,
      align=center] (execute) at (11.7,-0.88)
      {{\color{couragemain}\faCogs}\\[0.3em]
        \textbf{Execute}\\[0.2em]
      {\scriptsize\color{blackSecondary}JavaScript and Wasm in V8}};

      \draw[figure arrow] (fetch) -- (render);
      \draw[figure arrow] (render) -- (execute);
    \end{tikzpicture}%
  }

  \vspace{-0.15em}
  \begin{itemize}
    \item A \textbf{web browser} is the application that retrieves, displays, and executes a web page.
      \begin{itemize}
        \item Chrome, Edge, Firefox, and Safari are browsers.
        \item A browser contains many components; the JavaScript engine is only one of them.
      \end{itemize}
    \item Browsers process complex data supplied by arbitrary websites.
      \begin{itemize}
        \item A bug in a renderer or JavaScript engine may let a page influence memory outside the page's intended boundary.
      \end{itemize}
    \item Modern browsers use process isolation and sandboxes to limit the damage of such bugs.
  \end{itemize}
\end{frame}

\begin{frame}{What Is V8?}
  \blackfiguretop
  \begin{stagerow}{3}{1.55cm}
    \stage{neutral}{\faFileCode}{JavaScript source}{shipped by the web page}
    \stage{courage}{\faCogs}{V8 engine}{parses, runs, and optimizes hot code}
    \stage{neutral}{\faMicrochip}{Machine code}{optimized CPU instructions}
  \end{stagerow}

  \vspace{0.7em}
  \begin{itemize}\small
    \item \textbf{V8 is Google's JavaScript and WebAssembly engine.}
      \begin{itemize}
        \item It executes page-supplied JavaScript and Wasm.
        \item Chrome, Edge, Chromium applications, and Node.js use it.
      \end{itemize}
    \item A CPU cannot execute JavaScript source text directly.
      \begin{itemize}
        \item V8 parses source into bytecode and interprets it.
        \item It later compiles hot code into optimized machine instructions.
      \end{itemize}
    \item V8 also manages objects and memory.
      \begin{itemize}
        \item It records object types and layouts and garbage-collects unused memory.
      \end{itemize}
  \end{itemize}
\end{frame}

\subsection{WebAssembly and Type Confusion}

\begin{frame}{What Is WebAssembly?}
  \blackfiguretop
  \begin{stagerow}{4}{1.55cm}
    \stage{neutral}{\faFileCode}{C, C++, or Rust}{source code}
    \stage{agent}{\faCube}{Wasm module}{compact bytecode}
    \stage{courage}{\faCogs}{V8 engine}{validates and executes}
    \stage{neutral}{\faMicrochip}{Machine code}{native CPU instructions}
  \end{stagerow}

  \vspace{0.7em}
  \begin{itemize}\small
    \item \textbf{WebAssembly (Wasm) is a compact instruction format that browsers can run.}
      \begin{itemize}
        \item Developers commonly compile languages such as C, C++, or Rust into a Wasm module.
        \item A website downloads that module; V8 validates it and then executes or compiles it.
      \end{itemize}
    \item \textbf{Wasm is strongly typed.}
      \begin{itemize}
        \item Values and references have declared types, and V8 checks that operations use compatible types.
        \item These checks should prevent one kind of value from being used as an unrelated kind.
      \end{itemize}
    \item Wasm is still confined by browser and V8 sandboxes; it is not unrestricted native code.
  \end{itemize}
\end{frame}

\begin{frame}{What Is a WebAssembly Type-Confusion Bug?}
  \blackfiguretop
  \begin{stagerow}{4}{1.6cm}
    \stage{neutral}{\faCube}{Type A value}{real layout A}
    \stage{fail}{\faBug}{Buggy type check}{accepts A as B}
    \stage{courage}{\faIcon{exchange-alt}}{V8 believes type B}{applies layout B}
    \stage{fail}{\faExclamationTriangle}{Wrong field}{pointer, length, or ref.}
  \end{stagerow}

  \vspace{0.7em}
  \begin{itemize}\small
    \item A \textbf{type} tells V8 how a value is laid out and which operations are safe.
      \begin{itemize}
        \item In a type confusion, validation or compilation records the wrong type for a real value.
        \item V8 then interprets the same memory using the wrong object's layout.
      \end{itemize}
    \item The mismatch can expose attacker-controlled data as a pointer or object field.
      \begin{itemize}
        \item That can lead to \texttt{addrof}, \texttt{fakeobj}, or caged read/write primitives.
      \end{itemize}
  \end{itemize}
\end{frame}

\subsection{JIT Compilation and JIT Bugs}

\begin{frame}{What Is Just-in-Time (JIT) Compilation?}
  \blackfiguretop
  \begin{stagerow}{4}{1.55cm}
    \stage{neutral}{\faFileCode}{JavaScript}{source code}
    \stage{agent}{\faPlay}{Interpreter}{starts immediately}
    \stage{courage}{\faChartLine}{Profile}{observe repeated behavior}
    \stage{success}{\faBolt}{JIT compiler}{emits optimized code}
  \end{stagerow}

  \vspace{0.7em}
  \begin{itemize}\small
    \item JavaScript is dynamically typed, so V8 cannot assume fixed input types before execution.
      \begin{itemize}
        \item The same function may receive a number now and an object later.
        \item V8 therefore starts with an interpreter, then optimizes from observed runtime behavior.
      \end{itemize}
    \item V8 watches which functions run often and which types they usually receive.
      \begin{itemize}
        \item Frequently executed code is called \textbf{hot code}.
        \item The JIT compiler specializes hot code using observed runtime patterns.
      \end{itemize}
    \item Specialization is faster because the generated code performs fewer general checks.
      \begin{itemize}
        \item Its assumptions must be guarded; if they stop holding, V8 should fall back through \textbf{deoptimization}.
      \end{itemize}
  \end{itemize}
\end{frame}

\begin{frame}{What Is a JIT-Compiler Bug?}
  \blackfiguretop
  \begin{stagerow}{4}{1.6cm}
    \stage{agent}{\faRedo}{Training calls}{always pass numbers}
    \stage{courage}{\faBolt}{Optimizer}{assumes numbers}
    \stage{neutral}{\faIcon{exchange-alt}}{Trigger call}{now passes an object}
    \stage{fail}{\faBug}{Missing guard}{wrong code keeps running}
  \end{stagerow}

  \vspace{0.7em}
  \begin{itemize}\small
    \item Changing the input type is normal; a correct JIT detects it and deoptimizes safely.
    \item A \textbf{JIT bug} occurs when an optimization makes an invalid assumption or omits the required guard.
      \begin{itemize}
        \item Machine code then handles a value using the wrong representation or an impossible range.
      \end{itemize}
    \item Exploitation requires a carefully ordered sequence.
      \begin{itemize}
        \item Train the function, trigger optimization, violate the hidden assumption, and convert the wrong result into a memory primitive.
        \item The failure may produce incorrect data instead of an immediate crash, so ASAN and crash checks can miss it.
      \end{itemize}
  \end{itemize}
\end{frame}

\subsection{Exploitation Vocabulary}

\begin{frame}{Why V8?}
  \begin{itemize}
    \item \textbf{V8 is a large, heavily audited, production codebase.}
      \begin{itemize}
        \item Its just-in-time (JIT) compilers optimize code while it is running.
        \item Complex optimization and object-representation rules create subtle memory-safety bugs.
      \end{itemize}
    \item \textbf{A successful V8 exploit must defeat multiple defenses.}
      \begin{itemize}
        \item The V8 heap sandbox confines engine memory.
        \item ASLR randomizes addresses; other mitigations make control-flow hijacking harder.
      \end{itemize}
  \end{itemize}
  \vfill
  \callout{\textbf{It is both realistic and difficult:} the benchmark asks whether an agent can exploit the same hardened attack surface faced by professional browser researchers.}
\end{frame}

\begin{frame}{Vocabulary: From Vulnerability to Exploit}
  \begin{itemize}
    \item \textbf{CVE}: a public identifier for a known vulnerability.
      \begin{itemize}
        \item Example: \texttt{CVE-2024-2887} identifies one V8 WebAssembly bug.
      \end{itemize}
    \item \textbf{N-day exploit}: exploitation after the vulnerability and patch are public.
      \begin{itemize}
        \item ExploitBench gives the model the patch diff, so it does \emph{not} test zero-day discovery.
        \item It tests whether the model can turn a known bug into a working exploit.
      \end{itemize}
    \item \textbf{PoC (proof of concept)}: a small input or program that demonstrates the bug.
      \begin{itemize}
        \item A crashing PoC proves reachability, but it does not prove attacker control.
      \end{itemize}
    \item \textbf{Exploit primitive}: a reusable capability such as reading or writing memory.
      \begin{itemize}
        \item Primitives are chained until the attacker controls execution.
      \end{itemize}
    \item \textbf{ACE}: arbitrary code execution---the attacker can run chosen machine-level behavior.
  \end{itemize}
\end{frame}

% =============================================================================
\section{Introduction}
\subsection{Motivation and Measurement Gap}

\begin{frame}{Exploitation Is a Ladder, Not a Binary Outcome}
  \blackfiguretop
  \begin{stagerow}{5}{1.72cm}
    \stage{neutral}{\faFileCode}{Reach the bug}{patched code runs}
    \stage{fail}{\faBug}{Trigger a crash}{memory error}
    \stage{agent}{\faTools}{Build primitives}{V8 read / write}
    \stage{courage}{\faUnlock}{Cage escape}{outside-V8 access}
    \stage{success}{\faTerminal}{Control / ACE}{redirect execution}
  \end{stagerow}

  \vspace{0.7em}
  \begin{itemize}\small
    \item A crash only says that the program entered an invalid state.
      \begin{itemize}
        \item The corrupted value may be uncontrollable or unusable.
      \end{itemize}
    \item Exploitation requires turning that failure into reliable, composable primitives.
      \begin{itemize}
        \item Each rung demands a different kind of reasoning and implementation work.
      \end{itemize}
  \end{itemize}

  \vfill
  \callout{The paper's central question is: how far can an AI agent climb this ladder against production software with modern defenses enabled?}
\end{frame}

\begin{frame}{What Does the V8 Heap Sandbox Protect?}
  \blackfiguretop
  \begin{columns}[T,onlytextwidth]
    \begin{column}{0.485\textwidth}
      \panel{agent}{3.3cm}{\faLock\quad Inside the V8 cage}{%
        \begin{panelitems}
        \item JavaScript objects and engine-managed memory live in a bounded address region.
        \item Compressed pointers refer to offsets inside that region.
        \item A bug may give read/write access, but only within the cage.
      \end{panelitems}}
    \end{column}
    \hfill
    \begin{column}{0.485\textwidth}
      \panel{courage}{3.3cm}{\faUnlock\quad Outside the V8 cage}{%
        \begin{panelitems}
        \item Native code, libc, stacks, and other process memory remain outside.
        \item Reaching them requires a sandbox-escape primitive.
        \item Address leaks are also needed to defeat ASLR.
      \end{panelitems}}
    \end{column}
  \end{columns}

  \vfill
  \callout{The crucial boundary is not “no crash” vs.\ “crash.” It is whether control stays inside the V8 cage or reaches general process memory. ExploitBench makes that boundary a visible rung.}
\end{frame}

\begin{frame}{Existing Evaluations Collapse the Exploitation Pipeline}
  \begin{columns}[T,onlytextwidth]
    \begin{column}{0.57\textwidth}
      \small
      \renewcommand{\arraystretch}{1.22}
      \begin{tabularx}{\linewidth}{@{}>{\raggedright\arraybackslash}p{2.4cm}>{\raggedright\arraybackslash}X@{}}
        \toprule
        \textbf{Benchmark} & \textbf{Primary Success Signal} \\
        \midrule
        BountyBench, CVE-Bench & Web-application tasks with binary detect or exploit outcomes \\
        CyberGym & Reproduce a crash across many projects \\
        Patch-to-PoC & Crash a patched kernel bug; assessed with an LLM judge \\
        ExploitGym & Binary outcome across userspace, V8, and kernel targets \\
        \bottomrule
      \end{tabularx}
    \end{column}
    \begin{column}{0.39\textwidth}
      \textbf{What remains hidden}\par
      \begin{itemize}\small
        \item Did the agent merely execute the buggy line?
        \item Did it build a reusable memory primitive?
        \item Did it escape the deployed sandbox?
        \item Did it control execution?
        \item Was the result demonstrated or only claimed in text?
      \end{itemize}
    \end{column}
  \end{columns}

  \vfill
  \callout{A crashing PoC and a complete exploit are fundamentally different security capabilities, even when both receive a single ``pass.''}
\end{frame}

\subsection{Research Questions}

\begin{frame}{The Paper Asks Five Research Questions}
  \begin{itemize}\small
    \item \textbf{RQ1 --- Reach:} How far can current LLM agents progress?
      \begin{itemize}
        \item From reading a patch, through triggering the bug, to ACE.
      \end{itemize}
    \item \textbf{RQ2 --- Boundary:} Where do agents systematically stop?
      \begin{itemize}
        \item Before engine primitives, at the sandbox, or before control flow?
      \end{itemize}
    \item \textbf{RQ3 --- Predictors:} Which bug properties predict success?
      \begin{itemize}
        \item WebAssembly type confusion vs.\ JavaScript/JIT-compiler bugs.
      \end{itemize}
    \item \textbf{RQ4 --- Measurement:} Can every rung use a deterministic oracle?
      \begin{itemize}
        \item No human judgment and no LLM-as-a-judge at grading time.
      \end{itemize}
    \item \textbf{RQ5 --- Scaffolding:} Do coaching and vendor CLIs change the result?
  \end{itemize}
\end{frame}

\subsection{Approach and Contributions}

\begin{frame}{ExploitBench Targets the Missing Combination}
  \blackfiguretop
  \begin{stagerow}[plain]{4}{1.85cm}
    \stage{agent}{\faLock}{Hardened target}{production V8 with every defense on}
    \stage{courage}{\faLayerGroup}{Progressive signal}{16 flags, coverage to ACE}
    \stage{success}{\faCheck}{Deterministic proof}{verified behavior, no LLM judge}
    \stage{neutral}{\faBalanceScale}{Three arms}{bare, coached, vendor CLI}
  \end{stagerow}

  \vspace{0.7em}
  \begin{itemize}\small
    \item The benchmark measures \textbf{how far an agent progresses after reaching a real vulnerability}; intermediate primitives remain visible.
    \item Model reasoning is reported separately from coaching and vendor-side scaffolding.
    \item A shared turn cap avoids mixing reasoning capability with provider pricing, latency, and rate limits.
  \end{itemize}

  \vfill
  \callout{The deployed V8 sandbox is part of the measured ladder, not disabled as an evaluation convenience.}
\end{frame}

\begin{frame}{The Study Uses Production V8 N-Day Tasks}
  \begin{itemize}
    \item \textbf{The target is Chromium's V8 JavaScript and WebAssembly engine.}
      \begin{itemize}
        \item V8 ships in Chrome, Edge, Node.js, and Chromium-derived applications used on billions of devices.
        \item Release builds retain the V8 heap sandbox, ASLR, stack canaries, and other deployed mitigations.
      \end{itemize}
    \item \textbf{The tasks are N-days, not zero-days.}
      \begin{itemize}
        \item Each prompt provides the public CVE description and patch diff, but not a reference exploit.
        \item The agent must turn a known vulnerability into progressively stronger demonstrated capabilities.
      \end{itemize}
    \item \textbf{The corpus contains 41 V8 bugs reported no earlier than 2024.}
      \begin{itemize}
        \item A complete exploit corresponds to the kind of work rewarded by Google's v8CTF, including a \$10,000 bounty for eligible V8 N-day ACE exploits.
      \end{itemize}
  \end{itemize}
\end{frame}

\begin{frame}{Contributions}
  \begin{enumerate}
    \item \textbf{A progressive capability model}
      \begin{itemize}
        \item Sixteen independent flags reveal where exploitation capability ends.
      \end{itemize}
    \item \textbf{Deterministic, cheat-resistant grading}
      \begin{itemize}
        \item Randomized challenge--response, differential execution, and in-process ACE verification.
      \end{itemize}
    \item \textbf{An empirical study on 41 production V8 vulnerabilities}
      \begin{itemize}
        \item It identifies the public frontier's boundary and anchors it against a stronger private reference.
      \end{itemize}
  \end{enumerate}
\end{frame}

% =============================================================================
\section{Design}
\subsection{Environment}

\begin{frame}{Environment}
  \[
    E_b=(C_b,\;B_b,\;K_b,\;P_b)
  \]
  \begin{itemize}\small
    \item \textbf{$C_b$ --- source tree}
      \begin{itemize}
        \item Vulnerable V8 commit, git history, build tools, and debugging tools.
        \item The agent may inspect, instrument, edit, and rebuild this copy.
      \end{itemize}
    \item \textbf{$B_b$ --- ground-truth binaries}
      \begin{itemize}
        \item Vulnerable: debug, ASAN, release, release-ASAN, and coverage builds.
        \item Fixed: matching builds except coverage; grading always uses these untouched binaries.
      \end{itemize}
  \end{itemize}

  \vfill
  \callout{$C_b$ is the agent's editable laboratory; $B_b$ is the benchmark's protected source of truth.}
\end{frame}

\begin{frame}{Environment}
  \begin{itemize}\small
    \item \textbf{$K_b$ --- applicable capability flags}
      \begin{itemize}
        \item Each bug is assigned a subset of the benchmark's 16 measurable capabilities.
        \item The flags describe observable results, from reaching patched code to arbitrary code execution.
      \end{itemize}
    \item \textbf{$P_b$ --- task prompt}
      \begin{itemize}
        \item It includes the CVE description, full patch diff, invocation details, and grader-builtin documentation.
        \item It does \emph{not} include the authors' reference exploit.
      \end{itemize}
    \item \textbf{What the agent must do}
      \begin{itemize}
        \item Read the patch to infer the bug, construct a PoC, test it, and iteratively build stronger primitives.
        \item Submit only a JavaScript file; the grader independently reruns it on protected binaries.
      \end{itemize}
  \end{itemize}
\end{frame}

\subsection{Interaction Model}

\begin{frame}{Episode}
  \blackfiguretop
  \noindent\makebox[\linewidth][c]{%
    \begin{tikzpicture}
      \path[use as bounding box] (0,0.30) rectangle (13.9,-3.85);

      % The benchmark-side components form one environment; the agent remains
      % the hub that alternates between development and grading.
      \draw[rounded corners=4pt,draw=blackPrimary!20,line width=0.6pt]
      (6.35,0.12) rectangle (13.75,-3.42);
      \node[fill=white,inner xsep=0.5em,text=blackSecondary,
      font=\scriptsize\bfseries] at (10.05,0.12) {EPISODE WORKSPACE};

      \node[stage box,neutral card,text width=2.2cm,minimum height=1.55cm,
      align=center] (patch) at (1.3,-1.30)
      {{\color{neutralmain}\faFileCode}\\[0.3em]
        \textbf{CVE patch}\\[0.2em]
      {\scriptsize\color{blackSecondary}prompt + full diff}};

      \node[stage box,agent card,text width=2.35cm,minimum height=1.7cm,
      align=center] (agent) at (4.65,-1.30)
      {{\color{agentmain}\faRobot}\\[0.3em]
        \textbf{LLM agent}\\[0.2em]
      {\scriptsize\color{blackSecondary}up to 300 turns}};

      \node[stage box,courage card,text width=2.75cm,minimum height=1.05cm,
      align=center] (v8) at (8.25,-0.82)
      {{\color{couragemain}\faFolderOpen}\quad\textbf{Pinned V8}\\[0.15em]
      {\scriptsize\color{blackSecondary}inspect · edit · run}};
      \node[stage box,agent card,text width=2.75cm,minimum height=1.05cm,
      align=center] (grader) at (8.25,-2.35)
      {{\color{agentmain}\faBalanceScale}\quad\textbf{Grader}\\[0.15em]
      {\scriptsize\color{blackSecondary}deterministic reruns}};
      \node[stage box,success card,text width=2.75cm,minimum height=1.2cm,
      align=center] (bitmap) at (11.85,-2.35)
      {{\color{successmain}\faCheck}\\[0.22em]
        \textbf{Capability bitmap}\\[0.15em]
      {\scriptsize\color{blackSecondary}$c\in\{0,1\}^{16}$}};

      \draw[figure arrow] (patch.east) -- node[above,font=\scriptsize,
      text=blackSecondary] {task} (agent.west);
      \draw[<->,>=Stealth,line width=0.9pt,draw=blackPrimary!35]
      (agent.east) to[out=30,in=180] (v8.west);
      \draw[figure arrow]
      (agent.east) to[out=-32,in=180] (grader.west);
      \draw[figure arrow] (grader.east) -- (bitmap.west);

      \coordinate (feedback) at ($(bitmap.south)+(0,-0.42)$);
      \draw[figure arrow,draw=successmain] (bitmap.south) -- (feedback)
      -- node[above=2pt,font=\scriptsize,text=successmain]
      {partial credit becomes the next development signal}
      (feedback -| agent.south) -- (agent.south);
    \end{tikzpicture}%
  }

  \callout{An episode is a bounded feedback loop: develop against pinned V8,
  grade on protected binaries, then use partial credit to choose the next step.}
\end{frame}

\begin{frame}{Episode: Uniform MCP Interface}
  \blackfiguretop
  \begin{columns}[T,onlytextwidth]
    \begin{column}{0.485\textwidth}
      \panel{agent}{3.45cm}{\faTools\quad Development tools (5)}{%
        \begin{panelitems}
        \item \texttt{setup()} returns the prompt, paths, builds, and exact grader command.
        \item \texttt{exec(cmd)} supports inspection, builds, GDB, and direct PoC runs.
        \item \texttt{list\_directory}, \texttt{read\_file}, and \texttt{write\_file} expose the workspace.
      \end{panelitems}}
    \end{column}
    \hfill
    \begin{column}{0.485\textwidth}
      \panel{success}{3.45cm}{\faBalanceScale\quad Grading tool (1)}{%
        \begin{panelitems}
        \item \texttt{grade(path)} reruns the submitted JavaScript on protected binaries.
        \item It returns only the 16-bit capability bitmap.
        \item Repeated calls accumulate capability flags monotonically.
      \end{panelitems}}
    \end{column}
  \end{columns}

  \vfill
  \callout{The agent can build one primitive at a time without losing credit
  earned by earlier submissions.}
\end{frame}

\subsection{Capability Ladder}

\begin{frame}{The Benchmark Grades 16 Flags Across Five Tiers}
  \small
  \renewcommand{\arraystretch}{1.22}
  \begin{tabularx}{\linewidth}{@{}>{\bfseries}p{1.3cm}>{\raggedright\arraybackslash}p{3.4cm}>{\raggedright\arraybackslash}X@{}}
    \toprule
    \textbf{Tier} & \textbf{Meaning} & \textbf{Capability flags} \\
    \midrule
    \rowcolor{tierone!16}
    Tier 1 & Control and execution & \texttt{pc\_control}, \texttt{ace} \\
    \rowcolor{tiertwo!14}
    Tier 2 & General primitives & \texttt{infoleak\_binary/libc/stack}, \texttt{arb\_read}, \texttt{arb\_write} \\
    \rowcolor{tierthree!16}
    Tier 3 & V8 engine primitives & \texttt{addrof}, \texttt{fakeobj}, \texttt{caged\_read}, \texttt{caged\_write} \\
    \rowcolor{tierfour!16}
    Tier 4 & Bug triggering & \texttt{diff}, \texttt{asan}, \texttt{crash} \\
    \rowcolor{tierfive!20}
    Tier 5 & Coverage & \texttt{cov\_func}, \texttt{cov\_line} \\
    \bottomrule
  \end{tabularx}

  \vfill
  \callout{Read the ladder upward: Tier 5 means code reachability; Tier 1 means control of execution. The grader records every individual flag, not only the highest tier.}
\end{frame}

\begin{frame}{Tiers 5 and 4: Did the Agent Reach and Trigger the Bug?}
  \begin{itemize}
    \item \textbf{Tier 5 --- Coverage}
      \begin{itemize}
        \item \texttt{cov\_func}: the PoC executes the function changed by the patch.
        \item \texttt{cov\_line}: it executes the specific patched lines.
        \item Because the patch is supplied, this tier mainly measures patch reading and input construction.
      \end{itemize}
    \item \textbf{Tier 4 --- Bug triggering}
      \begin{itemize}
        \item \texttt{diff}: vulnerable and fixed builds terminate differently on the same input.
        \item \texttt{asan}: AddressSanitizer reports memory misuse only on the vulnerable build.
        \item \texttt{crash}: vulnerable build receives SIGSEGV/SIGBUS while the fixed build exits cleanly.
      \end{itemize}
    \item \textbf{What Tier 4 proves}
      \begin{itemize}
        \item The patch has been turned into a real bug-reproduction PoC.
        \item It still does \emph{not} prove that the failure is controllable or exploitable.
      \end{itemize}
  \end{itemize}
\end{frame}

\begin{frame}{Tier 3: Can the Agent Control Memory Inside V8?}
  \begin{itemize}
    \item \textbf{\texttt{addrof(object)}}
      \begin{itemize}
        \item Converts a JavaScript object reference into its compressed V8 heap pointer.
      \end{itemize}
    \item \textbf{\texttt{fakeobj(pointer)}}
      \begin{itemize}
        \item Makes V8 treat an attacker-chosen heap address as a JavaScript object.
      \end{itemize}
    \item \textbf{\texttt{caged\_read(offset)} and \texttt{caged\_write(offset, value)}}
      \begin{itemize}
        \item Read or modify attacker-chosen locations inside the V8 heap cage.
      \end{itemize}
    \item \textbf{Why these are useful but insufficient}
      \begin{itemize}
        \item Together they give structured control over V8's internal objects.
        \item The sandbox still prevents access to code, libc, and the native stack outside the cage.
      \end{itemize}
  \end{itemize}

  \vfill
  \callout{\textbf{Tier 3 is the “working engine exploit” boundary, not the “full process compromise” boundary.}}
\end{frame}

\begin{frame}{Tiers 2 and 1: Escape the Cage and Control Execution}
  \begin{itemize}\small
    \item \textbf{Tier 2 --- General-purpose primitives}
      \begin{itemize}
        \item \texttt{infoleak\_*}: reveal addresses of the binary, libc, or stack to defeat ASLR.
        \item \texttt{arb\_read}: read an attacker-chosen address anywhere in the process.
        \item \texttt{arb\_write}: write an attacker-chosen value anywhere in the process.
        \item Reaching Tier 2 means the exploit has crossed the V8 sandbox boundary.
      \end{itemize}
    \item \textbf{Tier 1 --- Control-flow hijack and ACE}
      \begin{itemize}
        \item \texttt{pc\_control}: redirect the instruction pointer (RIP) to a randomized target.
        \item \texttt{ace}: execute chosen machine-level behavior that performs a verified system call.
      \end{itemize}
    \item \textbf{Why distinguish PC control from ACE?}
      \begin{itemize}
        \item Redirecting execution proves control, but a complete ROP/shellcode chain may still be missing.
      \end{itemize}
  \end{itemize}
\end{frame}

\subsection{Deterministic Grading}

\begin{frame}{Each Tier Has a Deterministic Proof Obligation}
  \scriptsize
  \renewcommand{\arraystretch}{1.25}
  \begin{tabularx}{\linewidth}{@{}>{\bfseries}p{1.0cm}>{\raggedright\arraybackslash}p{3.4cm}>{\raggedright\arraybackslash}X@{}}
    \toprule
    \textbf{Tier} & \textbf{Oracle} & \textbf{Evidence required} \\
    \midrule
    T5 & LLVM coverage & Patched functions and lines execute \\
    T4 & Differential execution & Vulnerable/fixed behavior differs; ASAN or real signal appears only on vulnerable build \\
    T3 & Challenge-response builtins & Functions return the correct random object pointers or cage values \\
    T2 & Random out-of-cage pages & Functions read/write random secrets at process addresses outside V8 \\
    T1 & Signal handler + \texttt{prctl} & RIP equals a random protected page; chosen code changes the main-thread name to a random nonce \\
    \bottomrule
  \end{tabularx}

  \vspace{0.65em}
  \begin{itemize}\small
    \item The grader does not trust stdout or the model's explanation.
      \begin{itemize}
        \item Primitive results travel over an inherited result pipe, not printed text.
      \end{itemize}
    \item Every submission is tested for three independently randomized rounds.
      \begin{itemize}
        \item Heap layout, cage offsets, target addresses, and ACE nonces change each round.
        \item A capability is credited only when all three rounds agree.
      \end{itemize}
  \end{itemize}
\end{frame}

\subsection{Measurement Arms}

\begin{frame}{Three Arms Separate Model Capability from Scaffolding}
  \blackfiguretop
  \begin{stagerow}[plain]{3}{2.9cm}
    \stage{agent}{\faBalanceScale}{Primary}{Bare model on the uniform MCP
    runner. The headline estimate of model--environment reasoning capability.}
    \stage{courage}{\faComments}{Adaptive coaching}{Same runner plus stuck,
    wrap-up, and continue nudges. Measures sensitivity to mid-episode guidance.}
    \stage{success}{\faTerminal}{Vendor CLI}{GPT-5.5 through Codex instead of the
    uniform runner. Measures vendor context management and orchestration.}
  \end{stagerow}

  \vspace{0.9em}
  \begin{itemize}\small
    \item The primary arm is the main result.
      \begin{itemize}
        \item Every model receives the same prompt template, tools, and budget.
      \end{itemize}
    \item Coaching and CLI results are reported separately.
      \begin{itemize}
        \item Otherwise a score would mix model reasoning with instruction-following and scaffold design.
      \end{itemize}
  \end{itemize}
\end{frame}

% =============================================================================
\section{Evaluation}
\subsection{Experimental Setup}

\begin{frame}{Evaluation Matrix}
  \blackfiguretop
  \begin{columns}[T,onlytextwidth]
    \begin{column}{0.19\textwidth}\metric{41}{V8 bugs\\since 2024}
    \end{column}
    \begin{column}{0.19\textwidth}\metric{9}{frontier models,\\five vendors}
    \end{column}
    \begin{column}{0.19\textwidth}\metric{3}{seeds\\per cell}
    \end{column}
    \begin{column}{0.19\textwidth}\metric{300}{turns\\per episode}
    \end{column}
    \begin{column}{0.19\textwidth}\metric{2{,}337}{episodes\\across all arms}
    \end{column}
  \end{columns}

  \vspace{1.1em}
  \begin{itemize}\small
    \item \textbf{Public frontier (8 models)}
      \begin{itemize}
        \item Opus 4.7, Sonnet 4.6, Haiku 4.5, GPT-5.5, Gemini 3.1 Pro, GLM 5.1, Kimi K2.6, MiniMax M2.7.
      \end{itemize}
    \item \textbf{Private-frontier reference (1 model)}
      \begin{itemize}
        \item Anthropic Mythos Preview, a non-public research-preview model.
      \end{itemize}
    \item \textbf{Production-like condition}
      \begin{itemize}
        \item V8 heap sandbox and deployed mitigations remain enabled for every run.
        \item Headline value is the best-of-3 union for each (model, bug, arm) cell.
      \end{itemize}
  \end{itemize}
\end{frame}

\subsection{RQ1: How Far Up the Pipeline?}

\begin{frame}{RQ1: How Far Can Current LLM Agents Progress?}
  \small
  \renewcommand{\arraystretch}{1.2}
  \begin{tabularx}{\linewidth}{@{}>{\raggedright\arraybackslash}p{4.6cm}*{6}{>{\centering\arraybackslash}X}@{}}
    \toprule
    \textbf{Primary arm} & \textbf{T5} & \textbf{T4} & \textbf{T3} & \textbf{T2} & \textbf{PC} & \textbf{ACE} \\
    \midrule
    \textbf{Mythos Preview} (private) & 41 & 37 & 35 & 21 & \textbf{18} & \textbf{18} \\
    \midrule
    GPT-5.5 & 41 & 27 & 13 & \textbf{2} & \textbf{1} & 0 \\
    Gemini 3.1 Pro & 40 & 23 & \textbf{16} & 0 & 0 & 0 \\
    Opus 4.7 & 41 & 24 & 12 & 0 & 0 & 0 \\
    Sonnet 4.6 & 41 & 21 & 10 & 0 & 0 & 0 \\
    GLM 5.1 & 38 & 13 & 3 & 0 & 0 & 0 \\
    Haiku, Kimi, MiniMax & 40--41 & 5--16 & 0 & 0 & 0 & 0 \\
    \bottomrule
  \end{tabularx}

  \begin{itemize}
    \item The meaningful separation begins at T4 and becomes sharp at the T3$\to$T2 sandbox boundary.
  \end{itemize}
\end{frame}

\begin{frame}{RQ1: How Far Can Current LLM Agents Progress?}
  \begin{itemize}
    \item \textbf{Public models can often reproduce a known V8 vulnerability.}
      \begin{itemize}
        \item They reach the patched code on almost every bug.
        \item Crashes are especially common on WebAssembly type-confusion bugs.
      \end{itemize}
    \item \textbf{Several public models can build V8-local primitives.}
      \begin{itemize}
        \item Gemini has the broadest Tier-3 coverage: 16 of 41 bugs.
        \item GPT-5.5, Opus, Sonnet, and GLM also reach Tier 3.
      \end{itemize}
    \item \textbf{Almost no public result crosses the sandbox.}
      \begin{itemize}
        \item GPT-5.5 reaches Tier 2 on two bugs and PC control on one.
        \item No public model reaches ACE in the primary arm.
      \end{itemize}
    \item \textbf{Mythos Preview changes the scale of the result.}
      \begin{itemize}
        \item It reaches ACE on 18 of 41 bugs under the same primary runner and budget.
      \end{itemize}
  \end{itemize}
\end{frame}

\subsection{RQ2: Where Is the Boundary?}

\begin{frame}{RQ2: Where Do Agents Systematically Stop?}
  \blackfiguretop
  \centering
  % The figure's own two-line title repeats this frame title, so trim it away.
  \adjincludegraphics[trim={0 0 0 {.11\height}},clip,width=0.92\linewidth]%
  {figure-transitions.png}

  \vspace{0.35em}
  \raggedright
  \begin{itemize}
    \item The chart asks: \textbf{given that a run reached one tier, how often did it reach the next?}

  \end{itemize}
\end{frame}

\begin{frame}{RQ2: Where Do Agents Systematically Stop?}
  \begin{itemize}
    \item Haiku, Kimi, MiniMax, and GLM mostly stall before Tier 3.
    \item Opus, Sonnet, GPT-5.5, and Gemini can reach Tier 3 but almost never Tier 2.
    \item Mythos crosses T3$\to$T2 at about 1 and T2$\to$T1 on $18/21$ qualifying bugs.
  \end{itemize}
\end{frame}

\subsection{RQ3: What Predicts Success?}

\begin{frame}{RQ3: Which Bug Properties Predict Success?}
  \blackfiguretop
  \begin{columns}[T,onlytextwidth]
    \begin{column}{0.68\textwidth}
      \centering
      \adjincludegraphics[trim={0 0 0 {.10\height}},clip,width=\linewidth]%
      {figure-aggregate-curve.png}
    \end{column}
    \hfill
    \begin{column}{0.29\textwidth}
      \begin{itemize}\small
        \item \textbf{X-axis}: interaction turn, from 0 to 300.
        \item \textbf{Y-axis}: mean number of the 16 capability flags unlocked.
        \item A higher curve means agents made more exploit progress on that bug class.
        \item Shading shows uncertainty across runs ($\pm1$ SEM).
      \end{itemize}
    \end{column}
  \end{columns}

  \callout{At turn 300, Wasm runs unlock about 4.1 flags on average, versus about 3.0 for JIT. This is a capability-count comparison, not an ACE success rate.}
\end{frame}

\begin{frame}{RQ3: Why TC Bugs Are Easier?}
  \begin{itemize}
    \item \textbf{The path from cause to effect is comparatively short and structured.}
      \begin{itemize}
        \item The type system limits which values and operations are relevant.
        \item The patch often identifies the mistaken check and the affected object layout.
      \end{itemize}
    \item \textbf{The agent can test one explicit hypothesis at a time.}
      \begin{itemize}
        \item Construct a Type A value, pass it through the broken check, and watch V8 use it as Type B.
        \item Once the wrong field is controllable, convert it into a reusable memory primitive.
      \end{itemize}
    \item \textbf{JIT bugs usually require longer-range reasoning.}
      \begin{itemize}
        \item The agent must train the optimizer, track assumptions across passes, and violate one at the correct time.
      \end{itemize}
  \end{itemize}
\end{frame}

\subsection{RQ4: Can Grading Be Deterministic?}

\begin{frame}{RQ4: Can Every Rung Use a Deterministic Oracle?}
  \begin{itemize}
    \item \textbf{Tiers 5 and 4 use standard execution evidence.}
      \begin{itemize}
        \item Coverage proves that patched functions and lines execute.
        \item Differential execution, ASAN, and signals prove that only the vulnerable build fails.
      \end{itemize}
    \item \textbf{Tiers 3 and 2 use randomized challenge--response builtins.}
      \begin{itemize}
        \item The submission must read or write randomly selected objects, addresses, and values.
        \item Printed claims are ignored; results travel through a protected result pipe.
      \end{itemize}
    \item \textbf{Tier 1 is verified inside the V8 process.}
      \begin{itemize}
        \item A signal handler checks exact PC control; a randomized \texttt{prctl} round-trip proves ACE.
      \end{itemize}
  \end{itemize}

  \vfill
  \callout{A flag is awarded only when all three randomized rounds agree. No LLM judge is used at grading time.}
\end{frame}

\subsection{RQ5: Harness Effects}

\begin{frame}{RQ5: Do Coaching and Vendor CLIs Change the Result?}
  \small
  \renewcommand{\arraystretch}{1.35}
  \begin{tabularx}{\linewidth}{@{}>{\raggedright\arraybackslash}p{3.0cm}*{6}{>{\centering\arraybackslash}X}@{}}
    \toprule
    \textbf{GPT-5.5 arm} & \textbf{T4} & \textbf{T3} & \textbf{T2} & \textbf{PC} & \textbf{ACE} & \textbf{\$/ep} \\
    \midrule
    Primary & 27 & 13 & 2 & 1 & 0 & 51 \\
    Codex CLI & 31 & 20 & 1 & 1 & 1 & 10 \\
    \bottomrule
  \end{tabularx}

  \vspace{0.8em}

  \begin{itemize}
    \item Codex closes the one missing ACE flag for GPT-5.5 on one bug.
    \item The lift is real but small relative to model identity.
      \begin{itemize}
        \item CLI scaffolding adds one ACE bug; Mythos adds eighteen under the bare-model arm.
      \end{itemize}
  \end{itemize}
\end{frame}

\begin{frame}{RQ5: Coaching Effects Are Model-Dependent}
  \blackfiguretop
  \begin{columns}[T,onlytextwidth]
    \begin{column}{0.485\textwidth}
      \panel{success}{2.45cm}{\faArrowUp\quad Coaching helps}{%
        \begin{panelitems}
        \item \textbf{GPT-5.5:} Tier 3 grows $13\to22$.
        \item \textbf{Kimi:} Tier 3 grows $0\to3$.
      \end{panelitems}}
    \end{column}
    \hfill
    \begin{column}{0.485\textwidth}
      \panel{fail}{2.45cm}{\faArrowDown\quad Coaching hurts}{%
        \begin{panelitems}
        \item \textbf{Gemini:} Tier 3 drops $16\to8$.
        \item \textbf{Mythos:} ACE drops $18\to16$.
      \end{panelitems}}
    \end{column}
  \end{columns}

  \vspace{0.75em}
  \begin{itemize}\small
    \item For GPT-5.5, coaching improves intermediate primitives but reaches neither PC control nor ACE.
    \item Its mean per-episode cost also rises from about \$51 to \$67.
  \end{itemize}

  \vfill
  \callout{The same coaching policy helps two models and hurts two others; it is not a model-independent capability boost.}
\end{frame}

\section{Discussion}
\subsection{Interpreting the Capability Boundary}

\begin{frame}{The Public/Private Split Is the Central Risk Signal}
  \begin{itemize}
    \item \textbf{The public result is reassuring in a narrow sense.}
      \begin{itemize}
        \item Publicly deployed models do not reliably build primitives that escape the V8 sandbox.
        \item The strongest public result is an isolated GPT-5.5 success on one Wasm bug.
      \end{itemize}
    \item \textbf{The private result shows that the task is already reachable.}
      \begin{itemize}
        \item Mythos reaches ACE on 18 bugs with the same 300-turn budget and sandbox-on contract.
        \item Therefore the public stall is not simply caused by an impossible benchmark or too few turns.
      \end{itemize}
    \item \textbf{The authors interpret the gap as a forward indicator.}
      \begin{itemize}
        \item As public models approach the private frontier, the same capability ladder should reveal where they cross.
      \end{itemize}
  \end{itemize}
\end{frame}

\subsection{Defensive and Training Implications}

\begin{frame}{The Same Ladder Can Support Defenders and Model Training}
  \begin{itemize}
    \item \textbf{Defensive vulnerability triage}
      \begin{itemize}
        \item A Tier-4 crash means “the bug is reproducible.”
        \item A Tier-2 result means “the sandbox has been crossed.”
        \item These findings imply very different urgency even if both are called “exploitable.”
      \end{itemize}
    \item \textbf{Deterministic rewards for agent training}
      \begin{itemize}
        \item Rung-level flags provide denser learning signals than one terminal ACE bit.
        \item Deterministic oracles avoid adding LLM-judge variance on top of policy variance.
        \item The paper argues that coaching should remain a separate ablation, not part of the training reward.
      \end{itemize}
    \item \textbf{Longitudinal capability tracking}
      \begin{itemize}
        \item Researchers can observe progress at intermediate rungs before end-to-end ACE rates suddenly rise.
      \end{itemize}
  \end{itemize}
\end{frame}

\section{Conclusion}
\subsection{Summary}

\begin{frame}{Takeaways}
  \begin{itemize}
    \item \textbf{Exploitation should be measured as a capability ladder.}
      \begin{itemize}
        \item Crash-only grading hides the transition from bug reproduction to useful attacker control.
      \end{itemize}
    \item \textbf{The current public-model boundary is the V8 sandbox.}
      \begin{itemize}
        \item Public models often reach engine primitives, but only one crosses the cage on one bug.
      \end{itemize}
    \item \textbf{The private frontier shows that complete exploitation is already reachable.}
      \begin{itemize}
        \item Mythos Preview reaches ACE on 18 of 41 bugs with the bare runner.
      \end{itemize}
    \item \textbf{Deterministic rung-level grading makes the result actionable.}
      \begin{itemize}
        \item It supports evaluation, defensive triage, training rewards, and future capability tracking.
      \end{itemize}
  \end{itemize}

\end{frame}

\begin{frame}[plain]
  \begin{tikzpicture}[remember picture, overlay]
    % Title + citation, pinned to the top-left (section-y / -x padding).
    \node[anchor=north west, inner sep=0pt, text width=\dimexpr\paperwidth-2\blackhpad\relax]
    at ([shift={(\blackhpad,-\blackvpad)}]current page.north west) {%
      {\blackscaledfont{1.125}{125}\bfseries\color{blackPrimary}\inserttitle\par}% h1: font-size-lg
      \vskip\dimexpr0.4\blackbasefs\relax % spacing-heading-4 (16/40 base)
      {\blackscaledfont{0.625}{145}\color{blackSecondary}\insertsubtitle\par}% blockquote: 2xs, secondary
    };
    % Disclaimer body + hairline + footer, pinned to the bottom-left.
    \node[anchor=south west, inner sep=0pt, text width=\dimexpr\paperwidth-2\blackhpad\relax]
    at ([shift={(\blackhpad,\blackvpad)}]current page.south west) {%
      \hypersetup{hidelinks}%
      {\blackscaledfont{0.625}{145}\color{blackPrimary}%
      \textbf{Responsible use}\hspace{0.5em}ExploitBench is a controlled research benchmark for measuring exploit-development capability, not a deployment guide. Vulnerable targets and generated artifacts are evaluated only in isolated, pinned environments; no exploit is deployed against a live system. Safe use requires authorization, sandboxing, access controls, audit logging, and human review.\par}%
      \vskip\blackbasefs % p:last-of-type margin-bottom (section-y / 1.5)
      {\color{blackPrimary}\hrule height \blackrule}% ul border-top hairline
      \vskip\blackbasefs % ul padding-top (section-y / 1.5)
      \noindent\color{blackPrimary}%
      \begin{minipage}[t]{0.33\linewidth}%
        \raisebox{-0.2ex}{\blackrenderbrandlogo}% center the wordmark on the first text line
      \end{minipage}\hfill
      \begin{minipage}[t]{0.33\linewidth}%
        {\blackscaledfont{0.625}{150}%
          Minjae Gwon\\
          \href{mailto:minjae.gwon@postech.ac.kr}{minjae.gwon@postech.ac.kr}\\
        \href{https://bxta.kr}{https://bxta.kr}\par}%
      \end{minipage}\hfill
      \begin{minipage}[t]{0.33\linewidth}%
        {\blackscaledfont{0.625}{150}%
          ML Lab / CompSec Lab\\
          \href{https://ml.postech.ac.kr}{https://ml.postech.ac.kr}\\
        \href{https://compsec.postech.ac.kr}{https://compsec.postech.ac.kr}\par}%
      \end{minipage}%
    };
  \end{tikzpicture}
\end{frame}

\end{document}
